
%Some preparation: 
%\\
%\sout{\ppm{[I don't really understand this formulation yet:]}
%Recall, for a series of functions, $P=\{ P_n\}_{n\in \Z}$, the Chebyshev recurrence is} 
%\\
\bajmm{We say that a series of functions, $P=\{ P_n\}_{n\in \Z}$, is a ``Chebyshev series" if it satisfies the Chebyshev recursion:}
\beq \label{eq:Cheby0} 
P_{n+1}(x) = x P_{n}(x) +P_{n-1} (x)   ,
\eq  
(see e.g. \cite[\S6.3.3]{Martin91} and references therein; cf. also e.g. \cite{KarlinStudden}).
%\ppm{\sout{When a series $P$ obeys the recurrence it is a Chebyshev series.}}
This means of course that the series is 
determined in principle by any two successive values, $P_i$ and $P_{i+1}$. 
Classically these would be with $i=0$, as we will recall shortly. 
Two series $P,P'$ are `parallel' if $P_n = P'_{n+i}$ for some $i$. 

%\ppm{[some examples here? could have classical and one where $i$ below is not 1, say?]}

\mdef  \label{pr:cheb1}
Observe that if any two successive functions $P_i$ and $P_{i+1}$ are (integral) polynomial then all $P_n$ are. If in addition $\deg(P_{i+1})=\deg(P_i) + 1$ \ppmm{for some $i$,} then $\deg(P_{k+1})=\deg(P_{k}) + 1$ for all $k>i$ 
%\ppm{[-this appears to have different nature than first one, since only claimed to ascend.]}\bajm{[fixed]}
%\\
and further, if in addition
%\ppm{[-meaning the degree is ascending?...]}\bajm{[yes]}
$P_{i+1}$ is monic 
%\ppm{[...in which case is it enough that $P_{i+1}$ is monic?]}\bajm{[true]}
then so is $P_k$ for all $k>i$. 

\mdef We call a Chebyshev series of polynomials $P$ ``reduced" if there is no common non-trivial (\textit{i.e.} degree $>0$) polynomial factor 
%\sout{\ppm{[what ring?]}}
among all $P_i$ (equivalently, among any two successive %\sout{\ppm{non-zero [????]}\bajm{[no actually, we don't want non-zero here]}} \ppm{[I don't understand. if a term is zero then the whole next term is common, but only for those two.]}\bajm{[not quite, that term (the first non-zero one) will be a factor of the whole series (if $P_n=0$ then the cheb. rec. implies that $P_{n+2}(x)=x P_{n+1}(x)$) so the statement as is holds - i.e. the only reduced series with a 0 term somewhere is the classical one]}
terms). 
\\ 
Any polynomial series $P$ is a (polynomial) multiple %\ppm{[-is proportionality normally by a constant?]}
of a reduced one $P'$; $P'$ can be obtained from $P$ by removing a maximal common factor among two successive terms. 

%\ppm{[maybe reorganise a tad here. just for readability. R2 etc]}

\mdef  \label{de:Chebyab}
Let us write $\PP^{a,b} $ for the Chebyshev series with initial conditions 
$\PP^{a,b}_0 = a$ and $\PP^{a,b}_1 = b$.  

We may abuse notation slightly and, provided $i$ is fixed, write $\PP^{a,b} $ 
for the parallel series that starts with $\PP^{a,b}_i =a$. 

%\sout{Observe that if any two successive terms are polynomial then all are. 
%\\
%And if $\deg(b)=\deg(a)+1$ then $\deg(\PP^{a,b}_{k+1} = \deg(\PP^{a,b}_k)+1$ for all $k>i$. }
%\ppm{[its getting quite long. could delete some of the chatty notes-to-self? (nice though they are):]}
As we will see, explicit usage of the $\PP^{a,b}$ notation is often cumbersome in practice for general series (since $a$ may be a lengthy expression), and we will largely restrict the notation to abstract manipulations. Instead we use more specific notations for interesting cases, as in \S\ref{ss:onecup0}.  


\mdef 
The classical initial conditions are $P_0(x) =0$ and $P_1(x) =1$ giving then 
$P_2(x) =x $; $P_3(x) = x^2-1$; $P_4(x) = x(x^2 -2)$; and so on. In our $\PP^{a,b}$ notation this classical series is $\PP^{0,1}$. 

\mdef 
For convenience we set $P^U:=\PP^{0,1}$ since this series is related to the classical series of ``U-type" Chebyshev polynomials, $U$, which can be defined by the identity
\beq U_n(\cos(\th))\sin(\th)= \sin((n+1)\th),\label{eq:trig}\eq
by $P^U_n(x)=U_{n-1}(x/2)$. The polynomials $U_n$ and $P^U_n$ are well defined for $n<0$ with the above formulae still holding here; indeed, one can use these to determine
\beq P^U_{-n}(x)=-P^U_{n}(x),\label{eq:negcheb} \eq
for any $n\in \Z$. 

\mdef   \label{pr:Cbasis} Note that $\{ P^U_1(x),P^U_2(x),\dots \}$ is a $\C$-basis for the polynomial ring $\C[x]$.

\ignore{{ %%% all these promoted now to macros/macros-main.tex   %%%
%\newcommand{\ptarmigan}{Cheby condition}
\newcommand{\ptarmigan}{Cheby}
\newcommand{\ramping}{ramping}
\newcommand{\Ramping}{Ramping}
}} 
%I didn't see this command, I have replaced "Chebyshev condition" everywhere with "\bajm{X} property". Once we decide an appropriate name we can just globally replace \bajm{X} everywhere... this won't cause any issues
%LOL


\subsubsection{The \Ramping\ Property}

%\ppm{[-Ha, I don't think I understand killing the macro though. why would one replace a macro that can instantly and uniformly be changed to write anything you decide, with a passive symbol that will have to be edited across multiple instances!?]}\bajm{[you are correct I didn't see it, still a global replacement of ``bajm{X}" should do the job]}

Here we discuss a property that Chebyshev series may 
achieve eventually. 

%not quite sure how to proceed now. I guess I'll try again. 
\newcommand{\rampsfrom}{ramps from}   %%a ramp being ascending at, plausibly, 45 degrees.  :-) 
% OK you don't like this either. never mind. I've added the \ramping macro above.


\mdef\label{def:chebycond} 
Let $P$ be a 
%\sout{\ppm{reduced (hence polynomial) [????]}}\bajm{[I think it is better to let these be part of the definition not preconditions so I get more mileage by calling something X...] } 
Chebyshev series. For some integer $N$, we say that the series $P$ 
%\ppm{\em{\rampsfrom} $N$ [????]}\bajm{[decent name, but it would be much easier if we can just find a replacement for X since the rest of the paper can then be fixed by a global replacement]}
has the ``\ramping\ property 
%\ppm{[-any reason for this name?]}\bajm{[No (other than that the classical chebyshev series give an example) I just called it something to avoid wasting time choosing a name (which I have a tendency to do), happy for a better name :)]}
for $n\geq N$'' if it is polynomial and reduced, and if, furthermore, for $n\geq N$, $P_n$ is monic and $\deg(P_{n+1})=\deg(P_{n})+1$. 
%\ppm{[is reduced not an all n condition? and polynomial?]}\bajm{[it is, I am trying to say the last two conditions only apply for $n\geq N$]}
%\ppm{[as defined I think \ramping\ for N implies ramping for N+1 and so on. do you want to make N the minimum one? we could say:]  Series $P$  {\em \rampsfrom} $N$ if $N$ is minimal. [?]}\bajm{[No i actually don't because the polynomials we see later (the $P_n^{(\l)}$) appear to ramp for $n\geq l+3$ but I can't see a way to prove this generally - whereas we have a proof for $n\geq l+4$]}
Such a series is uniquely determined by $N$ and the ratio $P_{N+1}(x)/P_{N}(x)$, a rational function.




%\sout{Here the series is determined not just in principle, but a simple re-parametrisation gives it in a useful closed form, that we will recall shortly.  }
%\sout{In our $\PP^{a,b}$ notation this classical series is $\PP^{0,1}$. }

\mdef \label{pa:chebslv} 
Suppose that a series of polynomials $P$ has the \ramping\ property for $n\geq N$, and suppose that $d=\deg(P_N)$. 
%\sout{We claim  \ppm{[? - not sure what's going on here]} that one can determine} 
Then there exist unique coefficients $a_k, b_i, b'_j \in \C$, 
%for $k=-(d+1),-d,\dots,d,d+1$, $i=1,\dots, d+1$, and $j=1,\dots,d+2$
such that
\begin{align*}
    P_N(x)&=\sum_{k=-(d+1)}^{d+1} a_{k} P^U_{k}(x)=\sum_{i=1}^{d+1} b_i P^U_{i}(x),\\
    P_{N+1}(x)&=\sum_{k=-(d+1)}^{d+1} a_{k} P^U_{k+1}(x)=\sum_{j=1}^{d+2} b'_j P^U_{j}(x).\end{align*}
%\ppm{[-what is $b_i$? hark back to \ref{pr:Cbasis} maybe?]}\bajm{[okay?]}
The $a_k$ can be determined as follows: first solve for the coefficients $b_i$ and $b_j'$ - these are unique by \ref{pr:Cbasis}. Then by using \eqref{eq:negcheb} we obtain the following full rank linear system relating the $a_k$ to the $b_i$ and $b_j'$:
\begin{align*}
    b_k&=a_k-a_{-k}, & b'_{d+2}&=a_{d+1}, & b'_{d+1}&=a_{d}, &b'_{k'}&=a_{k'-1}-a_{-(k'+1)},
\end{align*}
for $k=1,\dots,d+1$ and $k'=1,\dots, d$. In particular, monicity of $P_N$ and $P_{N+1}$ imply that $a_{d+1}=1$ and $a_{-(d+1)}=0$.

Since the relation \eqref{eq:Cheby0} is linear, it thus follows that for any $l\in \Z$, we have
\beq 
P_{N+l}(x)=\sum_{k=-(d+1)}^{d+1} a_{k} P^U_{k+l}(x)=P^U_{d+1+l}(x)+\sum_{k=-d}^{d} a_{k} P^U_{k+l}(x). 
\eq




\ignore{
\mdef 
If we start with $P'_2(x)=2x$ and $P'_3(x) = (x-1)(x+2) \;$ 
(see e.g. \cite[\S5.2]{ShBr})
then  $P'_4(x) = x(x^2+x-4)$,
\beq   \label{eq:l=0+}
P'_5(x) = x^4 + x^3 -5x^2 -x +2, \qquad 
P'_6(x) =  x(x-1)(x^3 +2x^2 -4x -6)   \qquad 
\eq  
and so on. 
A better notation is perhaps $P'_n(x) = P^{(2)}_n(x)$, 
since this sequence is related to the $\lambda = (2)$ branch - as we will explain shortly. 
\\ \ppm{
Suppose now we start with $P''_2(x) = 2$ and $P''_3(x) = x-1$  
- a better notation here is $P''_n(x)  = P^{(1^2)}_n(x)$.
Then ... }

\mdef  \label{pa:Cheby1}
Here are tables of values for various initial conditions 
(we say `initial' conditions, but such Chebyshev sequences are determined by any two successive terms; other than for $P$, the starting point for $n$ values is not yet locked down):
\[
\begin{array}{c|c|c|ccc}
n\setminus  & P_n(x)        & P'_n =P^{(2)}_n      & P''_n = P^{(1^2)}_n \\
2      &  x                 & 2x               & 2 \\
3      & x^2 -1             & (x-1)(x+2)        & x-1  \\
4      & x(x^2-2)           &  x(x^2+x-4)      & (x+1)(x-2)  \\ 
5      & (x^2-x-1)(x^2+x-1) & x^4+x^3-5x^2 -x +2  & (x^3-x^2-3x+1)  \\
6       & x(x^2-1)(x^2-3)   & x(x-1)(x^3+2x^2-4x-6) & (x-1)(x^3-4x-2) 
\end{array}
\]
and continuing (with some initial conditions suggested by $l=1$, cf. \ref{pa:l1}):
\[ \hspace{-.861cm} 
\begin{array}{c|c|c|ccc}
n\setminus  & P^{(3)}_n             & P^{(21)}_n      &  P^{(1^3)}_n \\
1    &    &  3(x^2-1) & \\
2      &  3(x+2)             & x(x^2-4)               & 3 \\
3      &    x(x+4)            & (x^4-7x^2+3)         & x-2  \\
4      & (x+1)(x^2+3x-6)        &  x(x^2-1)(x^2-7)      & (x+1)(x-3)  \\ 
5      & x(x^3+4x^2-4x-10)        & (x^6-9x^4+14x^2-3)  & (x^3-2x^2-4x+2)  \\
6       & (x^5+4x^4-5x^3-14x^2+3x+6) & %\ppm{??(x^7-...x^6+...)} 
x \left(x^{6}-10 x^{4}+22 x^{2}-10\right)
& (x^4-2x^3-5x^2+4x+3) \\ 
7   &  \dots  & 
%\left(x -1\right) \left(x +1\right) \left(x^{6}-10 x^{4}+21 x^{2}-3\right)
\left(x^2 -1\right) \left(x^{6}-10 x^{4}+21 x^{2}-3\right)
& \dots 
\end{array}
\]
%\ppm{[-so far  unfinished! ]}



%\newcommand{\PP}{{\mathtt P}}  %% hook for Chebys 


\mdef 
In general there is a Chebyshev sequence for each `branch' for each $l$, 
hence for each $\lambda\vdash l+2$, 
thus a sequence $\PP^\lambda$ 
\ppm{defined specifically by ...
}

The first perspective on this is that the gram matrices on a branch are `related' by the Chebyshev recursion. This is true on the nose if we look at the slice through the Bratelli diagram corresponding to 1-cup modules $\Deltaaa{n}{(n-2,\lambda)}$ with $\lambda\in \{ (3), (1^3) \}$. 
... }




